\chapter{État modulaire bicouche et résolvante spectrale}
\label{rm:chap:estado-modular-resolvente}

\section{L’opérateur angulaire comme antécédent commun}

L’opérateur constitutif du \cref{rm:ch:vacuum} admet une formulation préalable dans laquelle l’intensité et l’orientation occupent des coordonnées distinctes d’un même objet. Soit \(R=R^*=R^{-1}\) une involution de multiplicités spectrales un et un, et soient
\begin{equation}
 P_+=\frac{I+R}{2},
 \qquad
 P_-=\frac{I-R}{2}
 \label{rm:eq:modular-projectors}
\end{equation}
ses projecteurs. Pour
\begin{equation}
 x=A_{\rm rad}>0,
 \qquad
 y=C^*_{\rm rad},
 \qquad
 |y|<x,
 \label{rm:eq:modular-domain}
\end{equation}
nous définissons le générateur angulaire et son exponentielle positive par
\begin{equation}
 K_{\angle}=xI+yR,
 \qquad
 T_{x,y}=\ee^{-K_{\angle}}
 =q_+P_++q_-P_-,
 \label{rm:eq:modular-angular-operator}
\end{equation}
où
\begin{equation}
 q_+=\ee^{-x-y},
 \qquad
 q_-=\ee^{-x+y}.
 \label{rm:eq:modular-qpm}
\end{equation}
La condition \(|y|<x\) place les deux valeurs propres dans \((0,1)\). Le déterminant et le quotient spectral séparent les deux coordonnées :
\begin{equation}
 D_A:=\det T_{x,y}=q_+q_-=\ee^{-2x},
 \qquad
 \varrho_{\rm or}:=\frac{q_-}{q_+}=\ee^{2y}.
 \label{rm:eq:modular-determinant-ratio}
\end{equation}
Le premier caractère est pair sous l’échange des feuillets ; le second enregistre leur orientation. La récupération est immédiate :
\begin{equation}
 x=-\frac12\log D_A,
 \qquad
 y=\frac12\log\varrho_{\rm or}.
 \label{rm:eq:modular-recover-xy}
\end{equation}

La provenance causale de \(x\) et de \(y\) reste liée à l’état \(\mathrm{APP}\to\TRIT\to\TPK\). En particulier,
\begin{equation}
 x=\frac{50\pi}{9}\,\mathfrak A_\beta,
 \qquad
 D_A=\exp\!\left(-\frac{100\pi}{9}\mathfrak A_\beta\right),
 \label{rm:eq:modular-alpha-determinant}
\end{equation}
où \(\mathfrak A_\beta=\alpha_{\HMT}\) désigne le caractère terminal du macro-objet dodécaphasique déjà construit. La coordonnée \(y\) provient ensuite du lecteur régional et conserve l’orientation que le déterminant élimine.

\begin{proposition}[Séparation canonique de l’intensité et de l’orientation]
\label{rm:prop:modular-separation}
L’application
\[
 (x,y)\longmapsto(D_A,\varrho_{\rm or})
 =(\ee^{-2x},\ee^{2y})
\]
est un difféomorphisme entre le domaine \(x>0\), \(|y|<x\), et le domaine
\[
 0<D_A<1,
 \qquad
 D_A<\varrho_{\rm or}<D_A^{-1}.
\]
L’involution des feuillets agit par \(
(D_A,\varrho_{\rm or})\mapsto(D_A,\varrho_{\rm or}^{-1})
\).
\end{proposition}

\begin{proof}
Les formules inverses sont \cref{rm:eq:modular-recover-xy}. L’inégalité \(|y|<x\) équivaut à \(|\log\varrho_{\rm or}|<-\log D_A\) et donne l’intervalle déclaré. Remplacer \(y\) par \(-y\) conserve \(D_A\) et inverse le quotient.
\end{proof}

\section{Normalisation de l’état bicouche}

La trace de l’opérateur angulaire est
\begin{equation}
 \Tr T_{x,y}=2\ee^{-x}\cosh y.
 \label{rm:eq:modular-trace-T}
\end{equation}
Sa normalisation élimine exactement l’intensité \(x\) et conserve l’orientation :
\begin{equation}
 \boxed{
 \rho_y:=\frac{T_{x,y}}{\Tr T_{x,y}}
 =\frac{\ee^{-yR}}{2\cosh y}
 =p_+(y)P_++p_-(y)P_- ,}
 \label{rm:eq:modular-rho}
\end{equation}
avec
\begin{equation}
 p_+(y)=\frac{\ee^{-y}}{2\cosh y},
 \qquad
 p_-(y)=\frac{\ee^{y}}{2\cosh y},
 \qquad
 p_++p_-=1.
 \label{rm:eq:modular-probabilities}
\end{equation}
La famille \(\rho_y\) est fidèle pour tout \(y\in\R\). Son observable d’orientation vérifie
\begin{equation}
 \Tr(\rho_yR)=-\tanh y,
 \qquad
 \rho_{-y}=R_{\rm sh}\rho_yR_{\rm sh}^*,
 \label{rm:eq:modular-orientation-expectation}
\end{equation}
où \(R_{\rm sh}\) représente tout unitaire permutant \(P_+\) et \(P_-\). Cette réalisation distingue l’involution spectrale \(R\), qui mesure le signe de feuillet, de l’unitaire \(R_{\rm sh}\), qui échange les feuillets.

\begin{definition}[Hamiltonien holographique modulaire local]
\label{rm:def:holographic-modular-hamiltonian} L’Hamiltonien modulaire de l’état bicouche est
\begin{equation}
 \boxed{
 K_y^{\rm mod}:=-\log\rho_y
 =\log(2\cosh y)I+yR.}
 \label{rm:eq:modular-Hamiltonian}
\end{equation}
La composante scalaire normalise l’état et la composante \(yR\) conserve la mémoire orientée.
\end{definition}

L’adjectif \emph{holographique} indique la provenance de \(R\), \(P_\pm\) et \(y\) à partir de l’état enrichi et de ses deux feuillets. L’égalité \cref{rm:eq:modular-Hamiltonian} est une identité de calcul fonctionnel ; sa promotion au rang de générateur d’une dynamique temporelle exige la spécification du flot physique correspondant.

\section{Entropie, parité et divergence d’orientation}

\begin{theorem}[Géométrie informationnelle exacte des deux feuillets]
\label{rm:thm:modular-information-geometry} L’entropie de von Neumann, la parité et l’entropie relative entre les deux orientations vérifient
\begin{align}
 S(\rho_y)
 &:=-\Tr(\rho_y\log\rho_y)
 =\log(2\cosh y)-y\tanh y,
 \label{rm:eq:modular-entropy}\\
 S(\rho_y)&=S(\rho_{-y}),
 \label{rm:eq:modular-entropy-even}\\
 \boxed{
 D(\rho_y\Vert\rho_{-y})
 =2y\tanh y.}
 \label{rm:eq:modular-relative-entropy}
\end{align}
La divergence est positive pour \(y\ne0\) et s’annule exactement dans le feuillet non orienté \(y=0\).
\end{theorem}

\begin{proof}
De \cref{rm:eq:modular-Hamiltonian} et de \(\Tr(\rho_yR)=-\tanh y\) résulte
\[
 S(\rho_y)=\Tr(\rho_yK_y^{\rm mod})
 =\log(2\cosh y)-y\tanh y.
\]
Les deux termes sont pairs. De plus,
\[
 \log\rho_y-\log\rho_{-y}=-2yR,
\]
et, par conséquent,
\[
 D(\rho_y\Vert\rho_{-y})
 =-2y\Tr(\rho_yR)=2y\tanh y.
\]
Le produit \(y\tanh y\) est positif pour \(y\ne0\).
\end{proof}

L’entropie quantifie la distribution paire des feuillets, tandis que la divergence quantifie la séparation informationnelle de leurs orientations. Les dérivées
\begin{equation}
 \frac{\dd}{\dd y}S(\rho_y)=-y\operatorname{sech}^2y,
 \qquad
 \frac{\dd^2}{\dd y^2}\log(2\cosh y)
 =\operatorname{sech}^2y
 \label{rm:eq:modular-entropy-derivatives}
\end{equation}
identifient \(S(\rho_0)=\log2\) comme maximum et fournissent l’information de Fisher de la famille exponentielle. Au voisinage du feuillet symétrique,
\begin{equation}
 D(\rho_y\Vert\rho_{-y})
 =2y^2-\frac23y^4+O(y^6).
 \label{rm:eq:modular-relative-expansion}
\end{equation}
La première contribution est quadratique : une faible mémoire orientée possède un coût informationnel du second ordre.

La distance de trace entre orientations s’obtient également sous forme fermée :
\begin{equation}
 \frac12\lVert\rho_y-\rho_{-y}\rVert_1=|\tanh y|.
 \label{rm:eq:modular-trace-distance}
\end{equation}
Comme \(|y|\geq|\tanh y|\), l’identité \cref{rm:eq:modular-relative-entropy} implique le contrôle
\begin{equation}
 D(\rho_y\Vert\rho_{-y})
 \geq 2\left(\frac12
 \lVert\rho_y-\rho_{-y}\rVert_1\right)^2,
 \label{rm:eq:modular-pinsker-exact-control}
\end{equation}
compatible avec l’inégalité de Pinsker. Ainsi, la quantité orientée possède simultanément une lecture spectrale, entropique et métrique.

\section{Résolvante normalisée du macro-objet angulaire}

Soit
\begin{equation}
 \kappa_\beta:=4\pi\mathfrak A_\beta>0.
 \label{rm:eq:modular-kappa-beta}
\end{equation}
Nous définissons la fonction matricielle méromorphe
\begin{equation}
 \boxed{
 \mathfrak M_\beta(z)
 :=\kappa_\beta^{-1}(zI-T_{x,y})^{-1},
 \qquad
 z\in\C\setminus\{q_+,q_-\}.}
 \label{rm:eq:modular-resolvent}
\end{equation}
Sa décomposition spectrale est
\begin{equation}
 \mathfrak M_\beta(z)
 =\frac1{\kappa_\beta}
 \left(\frac{P_+}{z-q_+}+\frac{P_-}{z-q_-}\right).
 \label{rm:eq:modular-resolvent-spectral}
\end{equation}
La résolvante réunit en un seul objet l’échelle de couplage, les pôles angulaires, les projecteurs de feuillet, le déterminant et tout le calcul fonctionnel ultérieur.

\begin{theorem}[Injectivité par les deux premiers moments]
\label{rm:thm:modular-resolvent-injective} L’application
\[
 (\mathfrak A_\beta,T_{x,y})
 \longmapsto\mathfrak M_\beta
\]
est injective pour \(\mathfrak A_\beta>0\). Plus précisément, si
\begin{align}
 B_0&:=\lim_{|z|\to\infty}z\mathfrak M_\beta(z),
 \label{rm:eq:modular-B0}\\
 B_1&:=\lim_{|z|\to\infty}z^2
 \left(\mathfrak M_\beta(z)-z^{-1}B_0\right),
 \label{rm:eq:modular-B1}
\end{align}
alors
\begin{equation}
 B_0=\kappa_\beta^{-1}I,
 \qquad
 \kappa_\beta=\tau_2(B_0)^{-1},
 \qquad
 T_{x,y}=\kappa_\beta B_1,
 \label{rm:eq:modular-reconstruct-kappa-T}
\end{equation}
où \(\tau_2=\Tr/2\). Par conséquent, le comportement asymptotique récupère \(\mathfrak A_\beta=\kappa_\beta/(4\pi)\) et le moment suivant récupère l’opérateur angulaire complet.
\end{theorem}

\begin{proof}
Pour \(|z|>\lVert T_{x,y}\rVert\), la série de Neumann donne
\begin{equation}
 \mathfrak M_\beta(z)
 =\frac1{\kappa_\beta}
 \sum_{n\geq0}\frac{T_{x,y}^{n}}{z^{n+1}}.
 \label{rm:eq:modular-Laurent}
\end{equation}
Les coefficients de \(z^{-1}\) et de \(z^{-2}\) sont précisément \(B_0\) et \(B_1\). Les égalités \cref{rm:eq:modular-reconstruct-kappa-T} résultent de la prise de la trace normalisée et de la multiplication par \(\kappa_\beta\). Deux paires possédant la même résolvante ont les mêmes coefficients et coïncident donc.
\end{proof}

Si \(q_+\ne q_-\), les données spectrales se récupèrent également à partir des pôles et des résidus :
\begin{equation}
 \operatorname*{Res}_{z=q_\pm}\mathfrak M_\beta(z)
 =\kappa_\beta^{-1}P_\pm,
 \qquad
 \det\mathfrak M_\beta(z)
 =\frac{\kappa_\beta^{-2}}
 {(z-q_+)(z-q_-)}.
 \label{rm:eq:modular-resolvent-residues}
\end{equation}
En particulier, le dénominateur scalaire contient \(q_++q_-\) et \(q_+q_-=D_A\), et
\begin{equation}
 P_+=\frac{T-q_-I}{q_+-q_-},
 \qquad
 P_-=\frac{T-q_+I}{q_--q_+}.
 \label{rm:eq:modular-projector-recovery}
\end{equation}
Le cas \(y=0\) produit un unique pôle simple de multiplicité spectrale deux : \(T=\ee^{-x}I\). La résolvante conserve alors l’intensité et l’opérateur scalaire, tandis que le choix des projecteurs est dépourvu de contenu intrinsèque parce que l’orientation est exactement nulle.

\begin{corollary}[Calcul fonctionnel à partir de la résolvante]
\label{rm:cor:modular-functional-calculus} Pour toute fonction \(f\) holomorphe au voisinage de \(\{q_+,q_-\}\),
\begin{equation}
 f(T)=\frac{\kappa_\beta}{2\pi\ii}
 \int_\Gamma f(z)\mathfrak M_\beta(z)\,\dd z,
 \label{rm:eq:modular-Dunford}
\end{equation}
où \(\Gamma\) entoure les deux pôles. En particulier, la réponse du vide
\[
 \mathcal V=(I-T^{90})(I-T^{120})^{-1}
\]
se reconstruit à partir de \(\mathfrak M_\beta\) sans incorporer de données scalaires supplémentaires.
\end{corollary}

\begin{proof}
L’identité est la formule de Dunford--Riesz appliquée à \((zI-T)^{-1}=\kappa_\beta\mathfrak M_\beta(z)\). La fonction \(f(z)=(1-z^{90})/(1-z^{120})\) est holomorphe au voisinage du spectre, puisque \(0<q_\pm<1\).
\end{proof}

\section{Compatibilité avec le raffinement nonadique}

L’information précédente se prolonge par la tour
\begin{equation}
 \mathcal A_m=M_2(\C)\otimes M_{9^m}(\C),
 \qquad
 \iota_m(a)=a\otimes I_9.
 \label{rm:eq:modular-UHF-tower}
\end{equation}
Soient
\begin{equation}
 T_m=T\otimes I_{9^m},
 \quad
 R_m=R\otimes I_{9^m},
 \quad
 \mathfrak M_{\beta,m}(z)
 =\kappa_\beta^{-1}(zI-T_m)^{-1},
 \label{rm:eq:modular-UHF-objects}
\end{equation}
et soit \(E_m=\operatorname{id}_{M_2\otimes M_{9^m}}\otimes\tau_9\) l’espérance traciale de \(\mathcal A_{m+1}\) vers \(\mathcal A_m\).

\begin{theorem}[Naturalité résolvante et modulaire]
\label{rm:thm:modular-UHF-naturality} Pour tout \(m\geq0\) et tout \(z\) de l’ensemble résolvant,
\begin{align}
 E_m(T_{m+1})&=T_m,
 \label{rm:eq:modular-UHF-T}\\
 E_m\!\left(\mathfrak M_{\beta,m+1}(z)\right)
 &=\mathfrak M_{\beta,m}(z),
 \label{rm:eq:modular-UHF-resolvent}\\
 E_m\!\left(f(T_{m+1})\right)&=f(T_m)
 \label{rm:eq:modular-UHF-functional}
\end{align}
pour toute fonction continue \(f\) sur le spectre. De même, les fonctionnelles d’état
\begin{equation}
 \omega_{y,m}(a)
 :=(\Tr_2\otimes\tau_{9^m})
 \bigl((\rho_y\otimes I_{9^m})a\bigr)
 \label{rm:eq:modular-UHF-state}
\end{equation}
sont compatibles avec les inclusions et définissent un état \(\omega_{y,\infty}\) sur la limite UHF.
\end{theorem}

\begin{proof}
Tous les objets de \cref{rm:eq:modular-UHF-objects} ont la forme \(a\otimes I_{9^m}\). L’espérance traciale élimine le dernier facteur identité et conserve \(a\). La même observation prouve \cref{rm:eq:modular-UHF-functional} ; la compatibilité des états s’obtient en évaluant \(\omega_{y,m+1}(a\otimes I_9)\) et en utilisant \(\tau_9(I_9)=1\).
\end{proof}

L’entropie locale relative à la trace \(\Tr_2\otimes\tau_{9^m}\) reste égale à \(S(\rho_y)\), et la divergence entre orientations reste égale à \(2y\tanh y\). Le raffinement ajoute de la résolution nonadique et conserve le contenu bicouche. De cette manière, la construction vit dans la tour UHF et dans sa clôture traciale, au lieu de dépendre d’une présentation fortuite par matrices \(2\times2\).

\section{Exactitude de la résolvante modulaire, obstructions de normalisation et séparation de la lecture thermodynamique}

\begin{proofstatusbox}
La décomposition de \(T\), l’état \(\rho_y\), l’Hamiltonien modulaire, l’entropie, la divergence relative, l’injectivité de la résolvante et la naturalité UHF sont des résultats exacts une fois \((\mathfrak A_\beta,C^*,R)\) fixés. La réunion de ces identités dans la résolvante \(\mathfrak M_\beta\) constitue une nouvelle formalisation de l’architecture préexistante de l’auteur. La lecture de \(D(\rho_y\Vert\rho_{-y})\) comme travail physique requiert une température, un protocole et une dynamique ; jusqu’à leur fixation, son statut est celui d’une mesure informationnelle exacte. Les réalisations au moyen de Barbero, de CKM ou de Higgs sont des lectures ultérieures et conservent leurs opérateurs d’entrelacement respectifs.
\end{proofstatusbox}

\begin{falsifierbox}
Le bloc est réfuté si la normalisation de \(T\) conserve une dépendance envers \(x\), si \(K_y^{\rm mod}\) ne donne pas \(\rho_y\) par exponentiation, si l’entropie relative diffère de \(2y\tanh y\), si deux paires distinctes \((\mathfrak A_\beta,T)\) possèdent la même résolvante, ou si une espérance \(E_m\) modifie le calcul fonctionnel. La récupération des projecteurs en \(y=0\) constitue en outre un contrôle négatif : tout sélecteur de feuillet survivant à cette dégénérescence aurait été introduit de l’extérieur.
\end{falsifierbox}
